\section*{Overview}

Binary search on answer is a modeling trick rather than a data structure. The goal is to turn ``find the best value''
into ``can I achieve at least / at most this value?'' Once that yes/no question is monotone, ordinary binary search
does the rest.

\section*{When to Use It}

Use it when:

\begin{itemize}
  \item the statement asks for a minimum feasible or maximum feasible value,
  \item a candidate answer can be checked faster than trying all answers directly,
  \item the feasibility condition changes only once across the ordered search space.
\end{itemize}

\section*{Core Idea}

Pick a candidate answer \(x\). Write a checker:
\[
\texttt{can(x)} = \text{``the problem is solvable with answer } x \text{''}.
\]
If the set of true values is an interval, binary search finds the boundary.

\section*{Key Insight}

The difficult part is usually not the search loop. It is simplifying the optimization target until the checker becomes
monotone. That checker often ends up being greedy, prefix-sum based, or a smaller DP.

\section*{Operations / Main Technique}

\begin{itemize}
  \item define the answer space,
  \item prove the feasibility predicate is monotone,
  \item derive a checker,
  \item binary search the first true or last true value.
\end{itemize}

\paragraph{Worked example.}
For partitioning an array into at most \(k\) segments while minimizing the maximum segment sum, the checker greedily
starts a new segment whenever adding the next value would exceed the candidate limit. If limit \(x\) works, any larger
limit works too.

\section*{Correctness Intuition}

Binary search itself is correct once the monotone predicate exists. The real proof burden is showing that the checker
answers feasibility correctly and that feasibility indeed moves in only one direction as the answer changes.

\section*{Complexity Analysis}

If the checker costs \(T\) and the answer space has size \(A\), the total cost is \(O(T \log A)\).

\section*{Implementation}

The code keeps a generic \texttt{first\_true} template and one sample checker for the segment-partition problem above.
That captures the pattern without tying it to one statement.

\section*{Common Pitfalls}

\begin{itemize}
  \item Searching before proving monotonicity.
  \item Choosing bounds that do not contain the answer.
  \item Writing a checker that is too slow, making the outer binary search irrelevant.
  \item Using integer search for a problem that really needs real-valued precision control.
\end{itemize}

\section*{Variants / Extensions}

\begin{itemize}
  \item Real-valued answer search.
  \item Parallel binary search for many offline queries.
  \item Parametric search with more structural checkers.
\end{itemize}

\section*{Practice Problems}

\begin{itemize}
  \item Minimize the largest segment sum.
  \item Maximize minimum distance under greedy placement.
  \item Any ``minimum feasible'' scheduling or capacity problem with a monotone checker.
\end{itemize}

\section*{References}

\begin{itemize}
  \item \href{https://cp-algorithms.com/num_methods/binary_search.html}{cp-algorithms: Binary Search on the Answer}.
  \item \href{https://usaco.guide/silver/binary-search}{USACO Guide: Binary Search}.
  \item \href{https://codeforces.com/catalog}{Codeforces Catalog}.
\end{itemize}
